% ថ្នាក់ទី ១១ (វិទ្យាសាស្ត្រ) · ជំពូកទី ៦ · មេរៀនទី ២ · សៀវភៅភាគ ២ ទំព័រ ៩២–១០១
\chapter{បន្សំ}

\begin{objBox}[វត្ថុបំណង]
\begin{enumerate}[1]
  \item ញែកភាពខុសគ្នារវាងចម្លាស់ និងបន្សំ
  \item កំណត់ចំនួនបន្សំ និងចំនួនចម្លាស់អាចបែងចែកបាន
  \item ពន្លាតទ្វេធាដឺក្រេទី $2, 3, 4$
  \item រកមេគុណទ្វេធាតាមត្រីកោណប៉ាស្កាល់
  \item រកមេគុណទ្វេធាតាមរូបមន្តបន្សំ
  \item រកមេគុណតួមួយនៃផលបូកក្នុងការពន្លាតទ្វេធា ។
\end{enumerate}
\end{objBox}

\section{បន្សំ}
\subsection{បន្សំនៃ \texorpdfstring{$r$}{r} ធាតុយកពី \texorpdfstring{$n$}{n} ធាតុ}

\begin{example}[ឧទាហរណ៍]
ក្នុងក្រុមមួយមានក្មេង $4$ នាក់ឈ្មោះ $A, B, C$ និង $D$ ។ បើក្មេង $2$ ចាប់ដៃសំដែងការគួរសមនឹងគ្នា តើការចាប់ដៃនេះមានប៉ុន្មានរបៀបខុសៗគ្នា?
\end{example}
\begin{itemize}
  \item បើគេឱ្យក្មេង $2$ នាក់អង្គុយជិតគ្នាដើម្បីថតរូប គេឃើញថារបៀបនៃការអង្គុយនេះមាន $2$ របៀបខុសគ្នា ព្រោះថាបើ $A$ អង្គុយខាងឆ្វេងនៃ $B$ ខុសពីរបៀបដែល $B$ អង្គុយខាងឆ្វេងនៃ $A$ ។ របៀបដែលឱ្យក្មេងម្ដង $2$ នាក់ក្នុងចំណោមក្មេង $4$ នាក់អង្គុយនេះជាចម្លាស់ $2$ ធាតុយកពី $4$ ធាតុ ហើយចំនួនចម្លាស់មាន ${}_4P_2 = \dfrac{4!}{(4 - 2)!} = \dfrac{4\times 3\times 2!}{2!} = 12$ របៀប ។
  \item បើគេឱ្យក្មេង $2$ នាក់ចាប់ដៃសំដែងការគួរសមនឹងគ្នា នោះគេឃើញរបៀបចាប់ដៃនេះមានតែមួយរបៀបព្រោះថា $A$ ចាប់ដៃ $B$ ក៏ដូចជា $B$ ចាប់ដៃ $A$ ដែរ មានន័យថាការចាប់ដៃនេះគ្មានលំដាប់មុនក្រោយ ឬឆ្វេងស្ដាំទេ ។
  \item ការដែលក្មេងចាប់ដៃគ្នាម្ដង $2$ នាក់ក្នុងចំណោម $4$ នាក់ដោយមិនគិតលំដាប់មុនក្រោយ ឬឆ្វេងស្ដាំនេះហៅថា \textbf{បន្សំនៃ $2$ ធាតុយកពី $4$ ធាតុ} ។
  \item គេឃើញថា ចំនួនរបៀបដែលក្មេងចាប់ដៃគ្នាម្ដង $2$ នាក់ក្នុងចំណោម $4$ នាក់មានតែ $6$ របៀបគឺតិចជាងចំនួនចម្លាស់នៃ $2$ ធាតុយកពី $4$ ធាតុចំនួន $2!$ ដង ព្រោះថាបន្សំមួយរបៀបគេអាចចម្លាស់បាន $2!$ របៀប ។ ដូច្នេះ ចំនួនបន្សំ $2$ ធាតុយកពី $4$ ធាតុស្មើនឹង $6$ ។ គេកំណត់សរសេរ ${}_4C_2 = 6$ ។
\end{itemize}
គេបាន $\dfrac{{}_4P_2}{2!} = {}_4C_2$ ឬ ${}_4C_2 = \dfrac{{}_4P_2}{2!} = \dfrac{4!}{(4 - 2)!2!} = \dfrac{4\times 3\times 2!}{2!\times 2} = 6$ ។
\[
\begin{array}{|cc|l|}\hline
\text{ចម្លាស់} & \text{(គិតលំដាប់)} & \text{បន្សំ (មិនគិតលំដាប់)}\\ \hline
AB & BA & AB \ (AB = BA)\\
AC & CA & AC\\
AD & DA & AD\\
BC & CB & BC\\
BD & DB & BD\\
CD & DC & CD\\ \hline
\end{array}
\]

\begin{formulaBox}[ជាទូទៅ]
បន្សំនៃ $r$ ធាតុយកពី $n$ ធាតុ គឺជាការយកព្រមគ្នាម្ដង $r$ ធាតុចេញពី $n$ ធាតុខុសៗគ្នាដោយមិនគិតពីលំដាប់នៃការយកចេញ ។ ចំពោះគ្រប់ $r \in \mathbb{N}$, $n \in \mathbb{N}$, $r \le n$ ចំនួនបន្សំនៃ $r$ ធាតុយកពី $n$ ធាតុគឺ
\[ {}_nC_r = \frac{{}_nP_r}{r!} = \frac{n!}{(n - r)!r!} \]
គេអាចសរសេរនិមិត្តសញ្ញាបន្សំដោយ $C(n, r)$, $C_n^r$ ឬ $\dbinom{n}{r}$ ។
\end{formulaBox}

\begin{example}[លំហាត់គំរូ ១]
គេមាន $5$ ចំណុចដែលគ្មានចំណុចបីណារត់ត្រង់គ្នាទេ ។ តើគេអាចសង់បានត្រីកោណចំនួនប៉ុន្មាន ដោយប្រើ $5$ ចំណុចនេះជាកំពូល ។
\end{example}
\ans
ដើម្បីសង់ត្រីកោណមួយ គេត្រូវប្រើបីចំណុចរត់មិនត្រង់គ្នា ដែលជាកំពូលត្រីកោណ ហើយលំដាប់នៃកំពូលមិនសំខាន់ទេ ។ ចំនួនត្រីកោណដែលសង់បាន ជាចំនួនបន្សំនៃ $3$ ធាតុយកពី $5$ ធាតុ ។ ចំនួនត្រីកោណទាំងអស់ស្មើនឹង ${}_5C_3 = \dfrac{5!}{(5 - 3)!3!} = \dfrac{5!}{2!3!} = 10$ ។

\begin{example}[លំហាត់គំរូ ២]
ក្នុងក្លិបមួយមានកីឡាករ $10$ នាក់ ។ គេបង្កើតក្រុមដែលមានសមាជិក $3$ នាក់ក្នុងចំណោមកីឡាករ $10$ នាក់ធ្វើជាដៃគូហ្វឹកហាត់ ។ តើគេអាចរើសកីឡាករ $3$ នាក់បានប៉ុន្មានក្រុម ក្នុងចំណោម $10$ នាក់?
\end{example}
\ans
ការរើសកីឡាករ $3$ នាក់ដើម្បីបង្កើតក្រុមតូច គឺជាបន្សំនៃ $3$ ធាតុយកពី $10$ ធាតុ ព្រោះគេមិនបានចាត់ជាប្រធាន អនុប្រធាន និងសមាជិកទេ ។ ចំនួនក្រុមទាំងអស់គឺ ${}_{10}C_3 = \dfrac{10!}{(10 - 3)!3!} = \dfrac{10!}{7!3!} = 120$ ក្រុម ។

\begin{pro}
\begin{enumerate}[1]
  \item ក្នុងទូនៃបណ្ណាល័យមួយមានសៀវភៅគណិតវិទ្យា $4$ ក្បាលខុសៗគ្នា និងរូបវិទ្យា $2$ ក្បាលខុសៗគ្នាដែរ ។ សុកត្រូវការអានសៀវភៅទាំងអស់ ប៉ុន្តែអាចខ្ចីសៀវភៅម្ដងបានតែ $3$ ក្បាលក្នុងចំណោមសៀវភៅនៅក្នុងទូនោះ ។ ក. តើសុកអាចរើសសៀវភៅ $3$ ក្បាលសម្រាប់ខ្ចីបានប៉ុន្មានរបៀប? \quad ខ. បើសុកចង់បានសៀវភៅគណិតវិទ្យា $2$ ក្បាល និងរូបវិទ្យា $1$ ក្បាល តើសុកអាចរើសបានសៀវភៅប៉ុន្មានរបៀប?
  \item ថ្នាក់រៀនមួយមានសិស្សប្រុស $12$ នាក់ និងសិស្សស្រី $8$ នាក់ ។ គ្រូចង់ដឹងសម្ថភាពសិស្ស $3$ នាក់ ដោយឱ្យដោះស្រាយលំហាត់មួយនៅក្ដារខៀន ។ ក. តើគ្រូអាចហៅសិស្សម្ដង $3$ នាក់មកក្ដារខៀនបានប៉ុន្មានរបៀប? \quad ខ. តើក្រុមសិស្ស $3$ នាក់មានប៉ុន្មានរបៀប បើគ្រូចង់ឱ្យក្រុមនេះមានសិស្សប្រុស $2$ នាក់ និងសិស្សស្រីម្នាក់?
\end{enumerate}
\ans
\begin{enumerate}[1]
  \item ក. ${}_6C_3 = 20$ \quad ខ. ${}_4C_2\times{}_2C_1 = 12$
  \item ក. ${}_{20}C_3 = 1140$ \quad ខ. ${}_{12}C_2\times{}_8C_1 = 528$ ។
\end{enumerate}
\end{pro}

\subsection{លក្ខណៈនៃបន្សំ}

\begin{example}[ឧទាហរណ៍]
${}_nC_{n-r} = {}_nC_r$ ព្រោះ ${}_nC_{n-r} = \dfrac{n!}{(n - n + r)!(n - r)!} = \dfrac{n!}{r!(n - r)!} = {}_nC_r$ ។ ${}_nC_n = \dfrac{n!}{(n - n)!n!} = \dfrac{1}{0!}$ ព្រោះ $0! = 1$ នោះ ${}_nC_n$ ជាបន្សំនៃ $n$ ធាតុយកពី $n$ ធាតុនោះការធ្វើបន្សំនេះបាន $1$ របៀប ។ ដូច្នេះ ${}_nC_n = 1$ ។
\end{example}

\begin{example}[លំហាត់គំរូ]
\begin{enumerate}[a]
  \item គណនា ${}_5C_3$ និង ${}_5C_2$ រួចសន្និដ្ឋាន ។
  \item គណនា ${}_nC_0$, ${}_nC_1$ ។
\end{enumerate}
\end{example}
\ans
\begin{enumerate}[a]
  \item ${}_5C_2 = \dfrac{5!}{(5 - 2)!2!} = \dfrac{5\times 4\times 3!}{3!2!} = 10$ ។ ${}_5C_3 = \dfrac{5!}{(5 - 3)!3!} = \dfrac{5\times 4\times 3!}{2!3!} = 10$ ។ ដូច្នេះ ${}_5C_3 = {}_5C_2$ ។
  \item ${}_nC_0 = \dfrac{n!}{(n - 0)!0!} = \dfrac{1}{0!} = 1$ ។ ${}_nC_1 = \dfrac{n!}{(n - 1)!1!} = \dfrac{n(n - 1)!}{(n - 1)!} = n$ ។
\end{enumerate}

\begin{pro}
បង្ហាញថា ${}_3C_0 + {}_3C_1 + {}_3C_2 + {}_3C_3 = 2^3$ ។
\ans
$1 + 3 + 3 + 1 = 8 = 2^3$ ។
\end{pro}

\subsection{ចម្លាស់បែងចែកបាន}

\begin{example}[ឧទាហរណ៍ ១]
គេរៀបសៀវភៅ $9$ ក្បាលដាក់ជាជួរក្នុង $9$ កន្លែង ដែលសៀវភៅទាំង $9$ ក្បាលនោះមានសៀវភៅគណិតវិទ្យា $6$ ក្បាលដូចគ្នា និងរូបវិទ្យា $3$ ក្បាលដូចគ្នា ។ តើគេអាចតម្រៀបសៀវភៅនេះបានប៉ុន្មានរបៀប?
\end{example}
\begin{itemize}
  \item ករណីសៀវភៅ $9$ ក្បាលខុសៗគ្នា នោះគេអាចតម្រៀបសៀវភៅបានចំនួន $9!$ របៀប ។
  \item ប៉ុន្តែមានសៀវភៅគណិតវិទ្យា $6$ ក្បាលដូចៗគ្នា នោះការរៀបសៀវភៅគណិតវិទ្យា $6$ ក្បាលខុសៗគ្នាចំនួន $6!$ របៀប ចាត់ទុកថាមានតែមួយករណីប៉ុណ្ណោះ ។
  \item ដូចគ្នាដែរចំពោះការរៀបសៀវភៅរូបវិទ្យា $3$ ក្បាលដូចគ្នា នោះចំនួនតម្រៀបសៀវភៅខុសៗគ្នា $3$ ក្បាលមានចំនួន $3!$ របៀប ចាត់ទុកថាមានតែ $1$ របៀបតែប៉ុណ្ណោះ ។
\end{itemize}
ដូច្នេះ ការតម្រៀបសៀវភៅ $9$ ក្បាលដែលមានសៀវភៅគណិតវិទ្យា $6$ ក្បាលដូចៗគ្នា និងរូបវិទ្យា $3$ ក្បាលដូចគ្នា កំណត់ដោយ $\dfrac{9!}{6!3!} = \dfrac{9\cdot 8\cdot 7\cdot 6!}{6!\cdot 3\cdot 2\cdot 1} = 3\times 4\times 7 = 84$ របៀប ។

\begin{example}[ឧទាហរណ៍ ២]
រកចំនួនចម្លាស់តួអក្សរទាំងអស់របស់ពាក្យ $BATTAMBANG$ ។
\end{example}
គេឃើញថា តួអក្សរទាំងអស់នៃពាក្យ $BATTAMBANG$ មាន $10$ តួ ដែលក្នុងនោះមានតួអក្សរ $A$ ចំនួន $3$ ដង អក្សរ $T$ ចំនួន $2$ ដង អក្សរ $B$ ចំនួន $2$ ដង និងអក្សរ $3$ តួទៀតខុសៗគ្នា ។
\begin{itemize}
  \item ដំបូងគេគិតថា គេតម្រៀប $A$ បីតួដូចគ្នាក្នុងចំណោម $10$ កន្លែង ។ គេបានចំនួនរបៀបដែលរៀប $A$ គឺ ${}_{10}C_3$ ព្រោះរើសយក $3$ ទីតាំងក្នុងចំណោម $10$ ទីតាំងដោយមិនគិតលំដាប់ ។
  \item ក្រោយពីរៀប $A$ រួចគេរៀប $T$ ពីរតួដូចគ្នាទៅក្នុង $7$ កន្លែងដែលនៅសល់ ។ គេបានចំនួនរបៀបដែលគេរៀប $T$ គឺ ${}_7C_2$ ។
  \item ក្រោយពីរៀប $A$ និង $T$ រួចហើយនៅសល់ $5$ កន្លែងសម្រាប់រៀប $B$ ពីរតួដូចគ្នា ។ គេបានចំនួនរបៀបដែលរៀប $B$ គឺ ${}_5C_2$ ។
  \item នៅសល់ $3$ កន្លែង សម្រាប់អក្សរ $M$ មួយ ៖ ${}_3C_1$ ។ នៅសល់ $2$ កន្លែង សម្រាប់អក្សរ $N$ មួយ ៖ ${}_2C_1$ ។ នៅសល់មួយកន្លែងសម្រាប់អក្សរ $G$ គឺមាន $1$ របៀប ។
\end{itemize}
ដូច្នេះ ចំនួនចម្លាស់តួអក្សរទាំង $10$ របស់ពាក្យ $BATTAMBANG$ គឺ
\[ \begin{aligned}
&{}_{10}C_3\times{}_7C_2\times{}_5C_2\times{}_3C_1\times{}_2C_1\times{}_1C_1\\
&= \frac{10!}{7!3!}\times\frac{7!}{5!2!}\times\frac{5!}{3!2!}\times\frac{3!}{2!1!}\times\frac{2!}{1!1!}\\
&= \frac{10!}{3!2!2!} = 151200
\end{aligned} \]

\begin{note}[សម្គាល់]
\begin{itemize}
  \item ការធ្វើចម្លាស់តួអក្សរទាំង $10$ នៃពាក្យ $BATTAMBANG$ នេះបើគេឧបមាថាអក្សរទាំង $10$ តួនេះខុសៗគ្នា នោះចំនួនចម្លាស់មាន $10!$ ។
  \item បើ $A$ មានបីតួខុសគ្នា ការតម្រៀប $A$ មានចំនួន $3!$ របៀប ។ តែដោយ $A$ មាន $3$ តួដូចគ្នា នោះចម្លាស់ $10!$ នេះមាន $3!$ ដងតិចជាង ។ ដូចគ្នានេះដែរចំពោះតួ $T$ ពីរតួ និងតួ $B$ ពីរតួ ។
\end{itemize}
\end{note}

\begin{formulaBox}[ជាទូទៅ]
ចម្លាស់ $n$ វត្ថុដែលក្នុងនោះមានវត្ថុប្រភេទទី $1$, ប្រភេទទី $2$, $\dots$, ប្រភេទទី $k$ ហើយវត្ថុប្រភេទនីមួយៗមានចំនួន $r_1, r_2, \dots, r_k$ រៀងគ្នា $(r_1 + r_2 + \dots + r_k = n)$ ហៅថា ចម្លាស់នៃវត្ថុដូចៗគ្នា ។ ចំនួនចម្លាស់នៃវត្ថុដូចៗគ្នា កំណត់ដោយ $\dfrac{n!}{r_1!r_2!\dots r_k!}$ ។
\end{formulaBox}

\begin{example}[លំហាត់គំរូ]
តើគេអាចចម្លាស់តួអក្សរនៃពាក្យ $MATHEMATIC$ បានប៉ុន្មានរបៀប?
\end{example}
\ans
ពាក្យនេះមាន $10$ តួ ដែល $M$, $A$, $T$ មាន $2$ តួម្នាក់ៗ ៖ $\dfrac{10!}{2!2!2!1!1!1!1!} = 453600$ ។

\begin{pro}
គេតម្រៀបទៀនពណ៌ក្រហម $3$ ដើម ទៀនពណ៌ស $2$ ដើម និងទៀនពណ៌លឿង $2$ ដើម ។ តើគេអាចតម្រៀបទៀនទាំងនេះបានប៉ុន្មានរបៀប?
\ans
$\dfrac{7!}{3!2!2!} = 210$ ។
\end{pro}

\section{ទ្រឹស្ដីបទទ្វេធា}
\subsection{ការពន្លាតទ្វេធានៃ \texorpdfstring{$(1 + a)^n$}{(1+a)^n}}

\begin{example}[ឧទាហរណ៍]
គេមានកន្សោម $(x + y)$, $(1 + y)$, $(1 - x)$, $(a + b)$, $(1 + a)$, $\dots$ កន្សោមនីមួយៗ ហៅថា \textbf{ទ្វេធា} គឺជាផលបូក ឬផលដករវាង $2$ តួ ។ គេសិក្សាលើការលើកជាស្វ័យគុណនៃទ្វេធា $(1 + a)$ ។ គេដឹងថា
\[ (1 + a)^2 = (1 + a)(1 + a) = 1 + 2a + a^2 \]
\[ (1 + a)^3 = (1 + a)(1 + 2a + a^2) = 1 + 3a + 3a^2 + a^3 \]
\[
\begin{aligned}
(1 + a)^4 &= (1 + a)(1 + 3a + 3a^2 + a^3)\\
&= 1 + 4a + 6a^2 + 4a^3 + a^4
\end{aligned}
\]
\end{example}
ក្នុងការពន្លាតទ្វេធានៃ $(1 + a)^n$ គេទទួលបានលទ្ធផលដូចខាងក្រោម ៖
\[
\begin{array}{|c|c|c|}\hline
\text{និទស្សន្ត} & \text{ការពន្លាតទ្វេធានៃ } (1 + a)^n & \text{ចំនួនតួ}\\ \hline
n = 0 & 1 & 1\\
n = 1 & 1 + a & 2\\
n = 2 & 1 + 2a + a^2 & 3\\
n = 3 & 1 + 3a + 3a^2 + a^3 & 4\\
n = 4 & 1 + 4a + 6a^2 + 4a^3 + a^4 & 5\\
n = 5 & 1 + 5a + 10a^2 + 10a^3 + 5a^4 + a^5 & 6\\ \hline
\end{array}
\]
តាមការពន្លាតនេះ គេសង្កេតឃើញថាលទ្ធផលនៃការពន្លាតទ្វេធានៃ $(1 + a)^n$ ជាផលបូកនៃ $(n + 1)$ ឯកធា អថេរ $a$ ដែលនិទស្សន្តរបស់ $a$ កើនតាមលំដាប់ពី $0$ ដល់ $n$ ហើយមេគុណរបស់ឯកធានីមួយៗហៅថា \textbf{មេគុណទ្វេធា} មានទំនាក់ទំនងគ្នាដូចខាងក្រោម ៖
\[
\begin{array}{c|ccccccccccccc}
n = 0 & &&&&&& 1 &&&&&&\\
n = 1 & &&&&& 1 && 1 &&&&&\\
n = 2 & &&&& 1 && 2 && 1 &&&&\\
n = 3 & &&& 1 && 3 && 3 && 1 &&&\\
n = 4 & && 1 && 4 && 6 && 4 && 1 &&\\
n = 5 & & 1 && 5 && 10 && 10 && 5 && 1 &\\
n = 6 & 1 && 6 && 15 && 20 && 15 && 6 && 1
\end{array}
\]
\begin{itemize}
  \item ការសរសេរទំនាក់ទំនងមេគុណទ្វេធាបន្តបន្ទាប់គ្នាមានរាងជាត្រីកោណនេះហៅថា \textbf{ត្រីកោណប៉ាស្កាល់ (Pascal's Triangle)} ។ ចំនួនដែលនៅខាងចុងជួរដេកនីមួយៗស្មើនឹង $1$ ហើយចំនួននៅខាងក្នុងត្រីកោណតាមជួរដេកនីមួយៗ បានពីការបូក $2$ ចំនួនជាប់គ្នារបស់ជួរដេកខាងលើ ។
  \item គេសង្កេតឃើញថា ការគណនា $(1 + a)^n$ បានដោយគុណ $(1 + a)$ នឹងខ្លួនឯងចំនួន $n$ កត្តា ។ ដូច្នេះ មេគុណទ្វេធា គឺជាការបន្សំ $a$ ក្នុងកត្តាផលគុណ ។
\end{itemize}
បើ $n = 4$ នោះការពន្លាតទ្វេធានៃ $(1 + a)^4$ គេបានមេគុណទ្វេធា $1, 4, 6, 4, 1$ ដែល ៖ $1$ ជាចំនួនបន្សំនៃ $0$ ធាតុរបស់ $a$ ក្នុង $4$ កត្តា ៖ ${}_4C_0 = 1$; $4$ ជាចំនួនបន្សំនៃ $1$ ធាតុរបស់ $a$ ក្នុង $4$ កត្តា ៖ ${}_4C_1 = 4$; $6$ ៖ ${}_4C_2 = 6$; $4$ ៖ ${}_4C_3 = 4$; $1$ ៖ ${}_4C_4 = 1$ ។ ដោយប្រើចំនួនបន្សំ ${}_nC_r$ គេបានត្រីកោណប៉ាស្កាល់
\[
\begin{array}{c|ccccccccccc}
n = 1 & &&&& {}_1C_0 && {}_1C_1 &&&&\\
n = 2 & &&& {}_2C_0 && {}_2C_1 && {}_2C_2 &&&\\
n = 3 & && {}_3C_0 && {}_3C_1 && {}_3C_2 && {}_3C_3 &&\\
n = 4 & & {}_4C_0 && {}_4C_1 && {}_4C_2 && {}_4C_3 && {}_4C_4 &\\
n = 5 & {}_5C_0 && {}_5C_1 && {}_5C_2 && {}_5C_3 && {}_5C_4 && {}_5C_5
\end{array}
\]

\begin{formulaBox}[ជាទូទៅ]
ក្នុងការពន្លាតទ្វេធានៃ $(1 + a)^n$ ដែល $n \in \mathbb{N}$ គេបាន
\[ \begin{aligned} (1 + a)^n = {}&{}_nC_0 + {}_nC_1\times a + {}_nC_2\times a^2 + \dots\\ &+ {}_nC_r\times a^r + \dots + {}_nC_n\times a^n \end{aligned} \]
\end{formulaBox}

\begin{example}[លំហាត់គំរូ ១]
ដោយប្រើរូបមន្តពន្លាតនៃ $(1 + a)^n$ ចូរពន្លាតទ្វេធានៃ ៖ \quad ក. $(1 + 2x)^5$ \quad ខ. $(1 - b)^5$ ។
\end{example}
\ans
\begin{enumerate}[a]
  \item តាង $a = 2x$ គេបាន
  \[ \begin{aligned}
(1 + 2x)^5 &= ({}_5C_0)(2x)^0\\
&\quad + ({}_5C_1)(2x)^1\\
&\quad + ({}_5C_2)(2x)^2\\
&\quad + ({}_5C_3)(2x)^3 + ({}_5C_4)(2x)^4 + ({}_5C_5)(2x)^5\\
&= 1 + 5(2x) + 10(2x)^2\\
&\quad + 10(2x)^3 + 5(2x)^4 + (2x)^5
\end{aligned} \]
  ដូច្នេះ $(1 + 2x)^5 = 1 + 10x + 40x^2 + 80x^3 + 80x^4 + 32x^5$ ។
  \item យក $a = -b$ គេបាន $(1 - b)^5 = 1 + 5(-b) + 10(-b)^2 + 10(-b)^3 + 5(-b)^4 + (-b)^5$ ។ ដូច្នេះ $(1 - b)^5 = 1 - 5b + 10b^2 - 10b^3 + 5b^4 - b^5$ ។
\end{enumerate}

\begin{example}[លំហាត់គំរូ ២]
សរសេរ $4$ តួដំបូងនៃការពន្លាតទ្វេធានៃ ៖ \quad ក. $(1 + 3x)^7$ \quad ខ. $(1 - 4x)^8$ ។
\end{example}
\ans
\begin{enumerate}[a]
  \item $(1 + 3x)^7 = 1 + ({}_7C_1)(3x) + ({}_7C_2)(3x)^2 + ({}_7C_3)(3x)^3 + \dots = 1 + 7(3x) + 21(9x^2) + 35(27x^3) + \dots$ ។ ដូច្នេះ $(1 + 3x)^7 = 1 + 21x + 189x^2 + 945x^3 + \dots$ ។
  \item $(1 - 4x)^8 = [1 + (-4x)]^8 = 1 + ({}_8C_1)(-4x) + ({}_8C_2)(-4x)^2 + ({}_8C_3)(-4x)^3 + \dots = 1 + 8(-4x) + 28(16x^2) + 56(-64x^3) + \dots$ ។ ដូច្នេះ $(1 - 4x)^8 = 1 - 32x + 448x^2 - 3584x^3 + \dots$ ។
\end{enumerate}

\begin{formulaBox}[ជាទូទៅ]
ក្នុងការពន្លាតទ្វេធានៃ $(1 + a)^n$ គេបានមេគុណនៃ $a^k$ ស្មើនឹង ${}_nC_k$ ។
\end{formulaBox}

\begin{pro}
\begin{enumerate}[1]
  \item ពន្លាតទ្វេធានៃ $\left(1 + \dfrac{2}{x}\right)^5$ ។
  \item ស្រាយបញ្ជាក់ថា ${}_nC_0 + {}_nC_1 + \dots + {}_nC_n = 2^n$ ។
\end{enumerate}
\ans
\begin{enumerate}[1]
  \item $1 + \dfrac{10}{x} + \dfrac{40}{x^2} + \dfrac{80}{x^3} + \dfrac{80}{x^4} + \dfrac{32}{x^5}$
  \item យក $a = 1$ ក្នុងរូបមន្ត $(1 + a)^n$ គេបាន $2^n = {}_nC_0 + {}_nC_1 + \dots + {}_nC_n$ ។
\end{enumerate}
\end{pro}

\subsection{ទ្រឹស្ដីបទទ្វេធា}

ចំពោះទ្វេធានៃ $(a + b)^n$ គេអាចសរសេរ $(a + b)^n = \left[a\left(1 + \dfrac{b}{a}\right)\right]^n = a^n\left(1 + \dfrac{b}{a}\right)^n$ ។ តាមរូបមន្តពន្លាតនៃកន្សោម $(1 + a)^n$ គេបាន
\[
\begin{aligned}
a^n\left(1 + \frac{b}{a}\right)^n &\\
&= a^n\left[1 + {}_nC_1\left(\frac{b}{a}\right) + {}_nC_2\left(\frac{b}{a}\right)^2 + \dots + {}_nC_n\left(\frac{b}{a}\right)^n\right]
\end{aligned}
\]
\[ (a + b)^n = a^n + {}_nC_1\, a^{n-1}b + {}_nC_2\, a^{n-2}b^2 + \dots + b^n \]

\begin{theorem}[ទ្រឹស្ដីបទទ្វេធា]
បើ $n \in \mathbb{N}$ នោះ
\[ \begin{aligned} (a + b)^n = {}&{}_nC_0\, a^n + {}_nC_1\, a^{n-1}b + {}_nC_2\, a^{n-2}b^2\\ &+ {}_nC_3\, a^{n-3}b^3 + \dots + {}_nC_n\, b^n \end{aligned} \]
\end{theorem}

\begin{note}[សម្គាល់]
ក្នុងការពន្លាតទ្វេធានៃ $(a + b)^n$ មេគុណនៃ $a^{n-k}b^k$ គឺ ${}_nC_k$ ។
\end{note}

\begin{example}[លំហាត់គំរូ ១]
ពន្លាតកន្សោម ៖ \quad ក. $(a + b)^3$ \quad ខ. $(2 + 3x)^3$ ។
\end{example}
\ans
\begin{enumerate}[a]
  \item $(a + b)^3 = a^3 + {}_3C_1\times a^2b + {}_3C_2\times ab^2 + b^3 = a^3 + 3a^2b + 3ab^2 + b^3$ ។
  \item តាង $a = 2$, $b = 3x$ គេបាន $(2 + 3x)^3 = 2^3 + 3(2)^2(3x) + 3(2)(3x)^2 + (3x)^3 = 8 + 36x + 54x^2 + 27x^3$ ។
\end{enumerate}

\begin{example}[លំហាត់គំរូ ២]
ពន្លាតទ្វេធានៃកន្សោមខាងក្រោមតាមស្វ័យគុណចុះនៃ $x$ ៖ \quad ក. $(2x - 3)^5$ \quad ខ. $\left(x + \dfrac{1}{x^2}\right)^6$ ។
\end{example}
\ans
\begin{enumerate}[a]
  \item \[ \begin{aligned}
(2x - 3)^5 &= [2x + (-3)]^5\\
&= (2x)^5 + ({}_5C_1)(2x)^4(-3)\\
&\quad + ({}_5C_2)(2x)^3(-3)^2\\
&\quad + ({}_5C_3)(2x)^2(-3)^3 + \dots\\
&= 32x^5 + 5(16x^4)(-3)\\
&\quad + 10(8x^3)(9)\\
&\quad + 10(4x^2)(-27) + \dots\\
&= 32x^5 - 240x^4 + 720x^3 - 1080x^2 + \dots
\end{aligned} \]
  \item \[ \begin{aligned} \left(x + \frac{1}{x^2}\right)^6 &= x^6 + {}_6C_1\, x^5\left(\frac{1}{x^2}\right) + {}_6C_2\, x^4\left(\frac{1}{x^2}\right)^2\\ &\quad + {}_6C_3\, x^3\left(\frac{1}{x^2}\right)^3 + \dots\\ &= x^6 + 6x^3 + 15 + \frac{20}{x^3} + \dots \end{aligned} \]
\end{enumerate}

\begin{example}[លំហាត់គំរូ ៣]
រកតម្លៃ $a, b, c$ បើគេដឹងថា $\left(a - \dfrac{1}{2}x\right)^6 = b - 96x + cx^2 + \dots$ ។
\end{example}
\ans
\[ \begin{aligned}
\left(a - \frac{1}{2}x\right)^6 &= \left[a + \left(-\frac{1}{2}x\right)\right]^6\\
&= a^6 + {}_6C_1\, a^5\left(-\frac{1}{2}x\right)\\
&\quad + {}_6C_2\, a^4\left(-\frac{1}{2}x\right)^2 + \dots\\
&= a^6 - 3a^5x + \frac{15}{4}a^4x^2 + \dots
\end{aligned} \]
ដោយ $\left(a - \dfrac{1}{2}x\right)^6 = b - 96x + cx^2 + \dots$ $\Rightarrow b - 96x + cx^2 = a^6 - 3a^5x + \dfrac{15}{4}a^4x^2$ ។ គេបាន $\begin{cases} b = a^6\\ 96 = 3a^5\\ c = \frac{15}{4}a^4 \end{cases}$ ។ $3a^5 = 96 \Rightarrow a^5 = 32 \Rightarrow a = 2$, $b = a^6 = 64$ និង $c = \dfrac{15}{4}\times 2^4 = 60$ ។ ដូច្នេះ $a = 2$, $b = 64$, $c = 60$ ។

\begin{pro}
សរសេរ $4$ តួដំបូងតាមស្វ័យគុណកើននៃ $x$ ក្នុងការពន្លាតទ្វេធានៃកន្សោម ៖ \quad ក. $(2 + 3x)^6$ \quad ខ. $\left(\dfrac{1}{2x} - 2x^2\right)^8$ ។
\ans
ក. $64 + 576x + 2160x^2 + 4320x^3 + \dots$ \quad ខ. $\dfrac{1}{256x^8} - \dfrac{1}{8x^5} + \dfrac{7}{4x^2} - 14x + \dots$ ។
\end{pro}

\section*{លំហាត់អនុវត្តន៍}
\begin{enumerate}[1]
  \item គណនា ៖ \quad ក. $\dfrac{{}_8C_5}{{}_8C_3}$ \quad ខ. $\dfrac{{}_{10}C_6}{{}_9C_5}$ \quad គ. $\dfrac{{}_{n+1}C_p}{{}_nC_p}$ ។
  \ans
  ក. $1$ \quad ខ. $\dfrac{210}{126} = \dfrac{5}{3}$ \quad គ. $\dfrac{n + 1}{n + 1 - p}$ ។
  \item ក្នុងក្រុមពិសេសមានសមាជិក $14$ នាក់ ។ គេបង្កើតគណៈប្រតិភូមានសមាជិក $5$ នាក់រើសចេញពីក្រុមពិសេសនេះ ។ តើគេអាចបង្កើតគណៈប្រតិភូបានប៉ុន្មានរបៀប?
  \ans
  ${}_{14}C_5 = 2002$ ។
  \item គេរើសមនុស្សចាស់ $4$ នាក់ និងក្មេង $3$ នាក់ចេញពីក្រុមមួយដែលមានមនុស្សចាស់ $8$ នាក់ និងក្មេង $5$ នាក់ ។ តើគេអាចរើសបានប៉ុន្មានរបៀប?
  \ans
  ${}_8C_4\times{}_5C_3 = 70\times 10 = 700$ ។
  \item គេរើសអក្សរ $5$ តួចេញពីអក្សរ $A, B, C, D, E, F, G, H, I$ និង $J$ ។
  \begin{enumerate}[a]
    \item តើគេអាចរើសបានប៉ុន្មានរបៀប?
    \item តើគេអាចរើសបានប៉ុន្មានរបៀប បើក្នុងចំណោម $5$ តួអក្សរនោះ ៖ មានព្យញ្ជនៈតែ $3$ គត់ \ / \ គ្មានស្រៈសោះ \ / \ និងយ៉ាងតិចមានស្រៈមួយ ។
  \end{enumerate}
  \ans
  ក. ${}_{10}C_5 = 252$ \quad ខ. (ស្រៈ $A, E, I$ និងព្យញ្ជនៈ $7$) ព្យញ្ជនៈតែ $3$ ៖ ${}_7C_3\times{}_3C_2 = 105$ ។ គ្មានស្រៈ ៖ ${}_7C_5 = 21$ ។ យ៉ាងតិចមានស្រៈមួយ ៖ $252 - 21 = 231$ ។
  \item តើគេអាចចែកសិស្សបានប៉ុន្មានរបៀប បើដឹងថាមាន ៖
  \begin{enumerate}[a]
    \item សិស្ស $5$ នាក់ត្រូវចែកជា $2$ ក្រុមដែលក្រុមមួយមាន $2$ នាក់ និងមួយទៀតមាន $3$ នាក់ ។
    \item សិស្ស $9$ នាក់ចែកជា $2$ ក្រុម ដែលមួយក្រុមមានសិស្ស $4$ នាក់ និង $5$ នាក់ ។
    \item សិស្ស $9$ នាក់ត្រូវចែកជា $3$ ក្រុម ដែលមួយក្រុមមានសិស្ស $2$ នាក់ $3$ នាក់និង $4$ នាក់ ។
  \end{enumerate}
  \ans
  ក. ${}_5C_2 = 10$ \quad ខ. ${}_9C_4 = 126$ \quad គ. $\dfrac{9!}{2!3!4!} = 1260$ ។
  \item សរសេរ $3$ តួដំបូងតាមស្វ័យគុណកើននៃ $x$ ក្នុងការពន្លាតទ្វេធានៃកន្សោមខាងក្រោម ៖ \quad ក. $(2 - ax)^6$ \quad ខ. $(3y + 2x)^5$ ។
  \ans
  ក. $64 - 192ax + 240a^2x^2 + \dots$ \quad ខ. $243y^5 + 810y^4x + 1080y^3x^2 + \dots$ ។
  \item បង្ហាញថា $(1 + \sqrt{x})^5 - (1 - \sqrt{x})^5 = 10\sqrt{x} + 20x\sqrt{x} + 2x^2\sqrt{x}$ ។ ទាញរកតម្លៃ $(1 + \sqrt{2})^5 - (1 - \sqrt{2})^5$ ។
  \ans
  តួដែលមានស្វ័យគុណគូនៃ $\sqrt{x}$ លុបគ្នា ៖ $2\left[5\sqrt{x} + 10(\sqrt{x})^3 + (\sqrt{x})^5\right] = 10\sqrt{x} + 20x\sqrt{x} + 2x^2\sqrt{x}$ ។ យក $x = 2$ ៖ $10\sqrt{2} + 40\sqrt{2} + 8\sqrt{2} = 58\sqrt{2}$ ។
  \item សរសេរ $4$ តួដំបូងតាមស្វ័យគុណកើននៃ $y$ ក្នុងការពន្លាតទ្វេធានៃកន្សោម ៖
  \begin{tasks}(4)
    \task $(1 + y)^{10}$
    \task $(1 - y)^{12}$
    \task $(1 + y^2)^9$
    \task $(1 - 3y^2)^7$
  \end{tasks}
  \ans
  ក. $1 + 10y + 45y^2 + 120y^3$ \quad ខ. $1 - 12y + 66y^2 - 220y^3$ \quad គ. $1 + 9y^2 + 36y^4 + 84y^6$ \quad ឃ. $1 - 21y^2 + 189y^4 - 945y^6$ ។
  \item គេពន្លាតទ្វេធានៃកន្សោមនីមួយៗខាងក្រោម ៖ \quad ក. ក្នុងលទ្ធផលពន្លាតកន្សោម $(1 - x^2)^{20}$ រកតួដែលមាន $x^6$ ។ \quad ខ. ក្នុងលទ្ធផលពន្លាតកន្សោម $(1 + 4x)^7$ រកតួដែលមាន $x^2$ ។
  \ans
  ក. ${}_{20}C_3(-x^2)^3 = -1140x^6$ \quad ខ. ${}_7C_2(4x)^2 = 336x^2$ ។
\end{enumerate}
